% !TeX document-id = {1c0b4298-276a-4038-8e08-c4a1f4846da7}
% !TeX TXS-program:bibliography = txs:///biber
\documentclass[10pt,a4paper]{article}
\usepackage[T1]{fontenc}
\usepackage{graphicx}
\usepackage{mathtools}
\usepackage{amssymb}
\usepackage{amsthm}
\usepackage{thmtools}
\usepackage{xcolor}
\usepackage{nameref}
\usepackage{hyperref}
\usepackage{color}
\usepackage{float}

\usepackage[
backend=biber,
style=authoryear,
natbib=true
]{biblatex}
\addbibresource{../bibliography.bib}

\title{EAE 6029}
\author{ Exercícios sobre Modelos Lineares}
\date{}
\begin{document}
	\maketitle
	No que segue, sempre suponha que as variáveis aleatória relevantes estão definidas no mesmo espaço de probabilidade.
 \paragraph{Exercício 1} Seja $Z$ uma variável aleatória discreta, que toma valores no conjunto $\{z_1,\ldots, z_n\}$; e $Y$ uma variável aleatória com segundo momento finito. Defina o vetor aleatório $X  = (\mathbf{1}_{\{z_2\}}(Z),\ldots \mathbf{1}_{\{z_n\}}(Z))$. Mostre que o melhor preditor linear de $Y$ como função de um intercepto e $X$ é da forma:
 
 $$\alpha_* + \delta_*'X$$
 com
 $$ \alpha_* =  \mathbb{E}[Y|Z=z_1]$$
$$\delta_j =   \mathbb{E}[Y|Z=z_{j+1}]-\mathbb{E}[Y|Z=z_1]$$
onde definimos o objeto $\mathbb{E}[Y|Z=z_j]$ como:
$$\mathbb{E}[Y|Z=z_j] \coloneqq \frac{\mathbb{E}[Y\mathbf{1}_{\{z_j\}}(Z)]}{\mathbb{P}[\{Z = z_j\}]}\, .$$. Conclua que o melhor preditor linear coincide com $\mathbb{E}[Y|Z]$.

\paragraph{Exercício 2} Seja $h^*$ a função de esperança condicional de uma variável aleatória $Y$ com segundo momento finito dado uma variável aleatória real $X$, i.e. $h^*(X) = \mathbb{E}[Y|X]$. Compute o melhor preditor linear de $Y$ como função de um intercepto e $X$, para os casos abaixo. Ilustre graficamente a derivada de $h^*$ comparando-a com a inclinação do melhor preditor linear, nos pontos do suporte\footnote{O suporte de uma variável aleatória $X$ é definido como o menor conjunto fechado $C$ tal que $\mathbb{P}[X \in C]=1$.} de $X$. Identifique as regiões onde o melhor preditor linear mais ``erra'' a inclinação de $h^*$.

\begin{itemize}
	\item[a] $h^*(x) = x^2 $ e $X\sim U[0,1]$.
	\item[b] $h^*(x) = \min\{x^3,100\}$ e $X \sim \operatorname{Exponecial}(1)$.
	\item[c] $h^*(x) = x^2$ e $X \sim \operatorname{\chi}^2(1)$.
\end{itemize}

 \paragraph{Exercício 3} Seja $A$ uma matriz simétrica positiva semidefinida:
 \begin{itemize}
 	\item[a] Mostre que os autovalores de $A$ sempre são não negativos. \text{Dica: } Se $v\neq 0$ é um autovetor de um autovalor $\lambda$, $Av=\lambda v$.
 	\item[b] Mostre que o menor autovalor de $A$ é zero se, e somente se, existe $v \neq 0$, $v'Av = 0$. \text{Dica: } para suficiência, utilize a decomposição espectral de uma matriz simétrica.
 \end{itemize} 
 
 \paragraph{Exercício 4} Seja $(\Omega, \Sigma,\mathbb{P})$ um espaço de probabilidade, em que cada elemento $\omega \in \Omega$ corresponde a uma firma num setor de interesse, e $\mathbb{P}[A]$ denota a proporção de firmas do ``tipo'' A . Considere o exemplo, visto em aula, em que a função de produção no setor é $h(z,l) = zl^{\alpha}$, $\alpha < 1$. Suponha que, no momento da decisão de produção, uma firma $\omega \in \Omega$ na população somente observe um sinal $\xi(\omega)$ de sua produtividade. A distribuição do sinal $\xi$ na população é denotada por $\mathbb{P}_\xi$. As firmas então escolhem $l$ otimizando:
 
 $$l(\omega) \in \operatorname{argmax}_{l \geq 0} \xi(\omega) l^{\alpha}- w l\, ,$$
onde $w$ denota o salário (em unidades do bem produzido pelo setor), aqui tratado como constante fixa.\footnote{Em outras palavras, a incerteza capturada pelas variáveis aleatórias só se deve à amostragem no sentido clássico, e não à incerteza econômica.} Após a escolha de demanda por trabalho, a produtividade de cada firma é revelada como $z(\omega) = \xi(\omega)\epsilon(\omega)$, onde o ``choque'' ou ``surpresa'' na produtividade  não é antecipável com base no sinal, i.e. $\xi$ é independente de $\epsilon$  ($\mathbb{P}_{(\xi,\epsilon)}  = \mathbb{P}_{\xi}\otimes \mathbb{P}_{\epsilon}$). A produção final das firmas é dada por $Y(\omega) = z(\omega) l(\omega)$.

\begin{itemize}
	\item[a] Encontre o melhor preditor linear da variável aleatória $\log(Y)$ como função de um intercepto e $\log(l)$. O coeficiente associado a $\log(l)$ coincide com $\alpha$? Por quê?
	\item[b] Suponha que, caso $Y(\omega) - wl(\omega)<0$, a firma $\omega$ não opera, e, nesse caso, a demanda por trabalho e quantidade produzida não são observadas. Calcule o melhor preditor linear nesse caso. Ele corresponde ao obtido anteriormente? Por quê?
\end{itemize}

	\paragraph{Exercício 5: Interpretação do melhor preditor linear num modelo causal heterogêneo.} Considere uma tripla de variáveis aleatórias $(Y,D,X)$, definidas num espaço de probabilidade  $(\Omega, \Sigma,\mathbb{P})$, em que a distribuição conjunta de $(Y,D,X)$ reflete a distribuição em uma população de um resultado escalar ($Y$) de interesse, uma causa binária ($D$) de $Y$ observável de interesse, e outras $k$ causas observáveis ($X$) de $Y$ sobre as quais não temos interesse direto. Nós definimos os resultados potenciais $Y(0)$ e $Y(1)$ correspondentes ao experimento hipótetico de se amostrar um elemento da população de interesse e fixar externamente sua causa $D$ em $0$ ($Y(0)$) ou $1$ ($Y(1)$), mantidas as demais causas de $Y$, observáveis ou não, em seus valores originais. Com base nessa definição do modelo causal, segue que: 

$$Y = Y(0) + D(Y(1)-Y(0))\, ,$$
e a variável aleatória $\tau \coloneqq Y(1)-Y(0)$ representa o efeito causal de  $D$ sobre uma unidade amostraao da população de interesse. 

No que segue, nós vamos supor que, \emph{condicionalmente a $X$}:

$$Y(0) \text{ é independente de } D\, ,$$
isto é, para quaisquer eventos $A \in \mathcal{B}(\mathbb{R})$ e $B \in 2^{\{0,1\}}$:
$$\mathbb{P}[\{Y(0) \in A\}\cap \{D \in B\}|X ]=\mathbb{P}[Y(0) \in A|X] \mathbb{P}[D \in B|X]$$

\begin{itemize}
	\item[a] Proveja uma interpretação para a hipótese de independência condicional apresentada acima.
	\item[b]  Mostre que, sob a hipótese de independência condicional apresentada acima, se  $\mathbb{P}[D=1|X](1-\mathbb{P}[D=1|X])>0$ quase certamente, e \textbf{ao menos uma} das duas condições abaixo for satisfeita:
	
	\begin{enumerate}
		\item[i] $\exists \psi \in \mathbb{R}^k$, $\mathbb{E}[Y(0)|X] = \psi'X$.
		\item[ii] $\exists \phi \in \mathbb{R}^k$, $\mathbb{E}[D|X] = \phi'X$.
	\end{enumerate}
	
	Então o coficiente $\beta^*$ de $D$ no melhor preditor linear de $Y$ em $D$ e $X$ (com intercepto caso ele esteja incluso em $X$), {definido como}:
	
	$$(\beta^*, \gamma^*) \in \operatorname{argmin}_{b \in \mathbb{R}, c \in \mathbb{R}^k}\mathbb{E}[(Y-Db - c'X)^2]$$
	
	corresponderá a:
	
	$$\beta^* =\frac{\mathbb{E}[ \omega(X) \tau(X)]}{\mathbb{E}[\omega(X)]}\, ,$$
	onde $\omega(X) = \mathbb{P}[D=1|X](1-\pi_*'X)$, com $\pi_* \in \operatorname{argmin}_{v \in \mathcal{R}^k}\mathbb{E}[(D-v'X)^2]$,  e $\tau(X)$ é a variável aleatória que denota o efeito médio de tratamento na subpopulação tratada definida pelos valores de $X$, i.e.
	
	$$\tau(X) \overset{\text{def}}{=} \mathbb{E}[(Y(1)-(0))|D=1,X] = \frac{\mathbb{E}[(Y(1)-Y(0))D|X]}{\mathbb{P}[D=1|X]}$$
	
	\textit{Dica:} use as condições de primeira ordem do melhor preditor linear para mostrar que:
	$$\beta^* = \frac{\mathbb{E}[(D-\pi_*'X)Y]}{\mathbb{E}[(D-\pi_*'X)^2]}\, ,$$
	onde você deve mostrar que $\mathbb{P}[D=1|X](1-\mathbb{P}[D=1|X]) > 0$ q.c. implica que $\mathbb{E}[(D-\pi_*'X)^2] > 0$. Para concluir, você precisará mostrar que, sob a hipótese de independência condicional, se ao menos (i) ou (ii) vale, então:
	
	$$\mathbb{E}[(D-\pi_*'X)Y(0)] = 0 \, .$$
	
	\item[c] Usando o resultado anterior, mostrar que:
	
	$$\beta^* = \int \omega^*(x) \tau(x) \mathbb{P}_X(dx)\, ,$$
	onde $\mathbb{P}_X$ denota a lei de probabilidade induzida por $X$, e onde os \textbf{pesos} $\omega^*$ são tais que $\int \omega^*(x) \mathbb{P}_X(dx)  =  1$, . 	Quando a condição (ii) é satisfeita,  como esses pesos se comparam àqueles dados pelo efeito médio do tratamento na população tratada, definido como:
	
	\begin{equation*}
		\begin{aligned}					\tau^{\text{ATT}}\overset{\text{def}}{=}\mathbb{E}[(Y(1)-Y(0))|D=1] = \frac{\mathbb{E}[(Y(1)-Y(0))D]}{\mathbb{P}[D=1]} = \\ \int \tau(x) \frac{\mathbb{P}[D=1|X=x]}{\mathbb{P}[D=1]} \mathbb{P}_X(dx)\, ,
		\end{aligned}
	\end{equation*}
	
	Proveja uma intuição para esse resultado.
	
	\item[d] Proveja uma condição suficiente para que os pesos $\omega^*$ sejam não negativos (quase certamente). Mostre que os pesos não serão positivos quase certamente se existir uma entrada $X_j$ de $X$ que toma, para qualquer intervalo da reta, valores nesse intervalo com probabilidade positiva, e tal que $\pi_{*,j}\neq 0$. Mostre que os  $\omega^*$ serão não negativos se $X = (\mathbf{1}_{\{z_1\}}(Z), \mathbf{1}_{\{z_2\}}(Z),\ldots \mathbf{1}_{\{z_k\}}(Z))'$ onde $Z$ é uma variável aleatória que toma valores no conjunto $\{z_1,\ldots, z_k\}$.
\end{itemize}



 

\end{document}